% requer booktabs. Diagnostico descritivo.
\begin{table}[htbp]\centering
\caption{Administrative completeness markers, by year of first application}
\label{tab:D15_C_marcadores}
\small
\begin{tabular}{lrrrrrr}
\toprule
Year & N & CPF & Sex & Postcode & Race decl. & Priority \\
\midrule
2010 & 81,004 & 94.2 & 100.0 & 97.2 & 63.9 & 3.7 \\
2011 & 124,240 & 96.3 & 100.0 & 97.9 & 61.9 & 7.7 \\
2012 & 220,922 & 98.1 & 100.0 & 98.3 & 60.8 & 11.1 \\
2013 & 191,149 & 99.3 & 100.0 & 98.4 & 59.8 & 16.6 \\
2014 & 174,676 & 99.5 & 100.0 & 98.5 & 56.5 & 18.9 \\
2015 & 179,869 & 99.6 & 100.0 & 98.5 & 57.9 & 19.4 \\
2016 & 172,341 & 99.6 & 100.0 & 98.6 & 55.6 & 20.4 \\
2017 & 163,922 & 99.7 & 100.0 & 98.6 & 54.0 & 13.9 \\
2018 & 165,331 & 99.5 & 100.0 & 98.6 & 65.1 & 15.0 \\
\bottomrule
\end{tabular}
\\[3pt]{\footnotesize \textit{Note:} Share of first applications with each field present: taxpayer id of the guardian (CPF); sex; postcode matched to a census tract; race declared (categories 1--5, excluding not-declared); a priority code. Rising completeness would mark a change of form or registration system. Source: build04\_diag15.py.}
\end{table}
